\section{Resultados}

A Figura~\ref{fig:comparacao} compara visualmente a imagem original com as
versões em grayscale e segmentada. Os arquivos usados nesta comparação são
as versões renomeadas do triângulo em \texttt{report/images}.

\begin{figure}[H]
    \centering
    \begin{minipage}{0.31\textwidth}
        \centering
        \includegraphics[width=\linewidth,height=4.5cm,keepaspectratio]{original_triangulo.jpg}

        \small Imagem original
    \end{minipage}
    \hfill
    \begin{minipage}{0.31\textwidth}
        \centering
        \includegraphics[width=\linewidth,height=4.5cm,keepaspectratio]{grayscale_triangulo.jpg}

        \small Imagem em grayscale
    \end{minipage}
    \hfill
    \begin{minipage}{0.31\textwidth}
        \centering
        
\includegraphics[width=\linewidth,height=4.5cm,keepaspectratio]{segmented_triangulo.jpg}

        \small Imagem segmentada
    \end{minipage}
    \caption{Comparação entre a imagem original, a conversão para grayscale e a segmentação do triângulo.}
    \label{fig:comparacao}
\end{figure}

Inclua nesta seção as observações sobre as diferenças entre as imagens e,
se aplicável, tabelas com métricas quantitativas dos resultados.

\subsection{Exemplos positivos e negativos}

Para complementar a comparação anterior, foram selecionadas imagens que
representam exemplos positivos e negativos para a análise de textura. Os
arquivos foram copiados para \texttt{report/images} com nomes descritivos.

\subsubsection{Exemplos positivos}

\begin{figure}[H]
    \centering
    \triplaimagem
        {original_arvores-ceu.jpg}
        {grayscale_arvores-ceu.jpg}
        {segmented_arvores-ceu.jpg}
        {Árvores e céu}
    \triplaimagem
        {original_muro-ceu.jpg}
        {grayscale_muro-ceu.jpg}
        {segmented_muro-ceu.jpg}
        {Muro e céu}
    \triplaimagem
        {original_muro.jpg}
        {grayscale_muro.jpg}
        {segmented_muro.jpg}
        {Muro}
    \triplaimagem
        {original_floresta.jpg}
        {grayscale_floresta.jpg}
        {segmented_floresta.jpg}
        {Floresta}
    \triplaimagem
        {original_poste-grama.jpg}
        {grayscale_poste-grama.jpg}
        {segmented_poste-grama.jpg}
        {Poste e grama}
    \caption{Exemplos positivos e suas versões em grayscale e segmentadas.}
    \label{fig:exemplos-positivos}
\end{figure}

\subsubsection{Exemplos negativos}

\begin{figure}[H]
    \centering
    \triplaimagem
        {original_meio-fio.jpg}
        {grayscale_meio-fio.jpg}
        {segmented_meio-fio.jpg}
        {Meio-fio}
    \triplaimagem
        {original_rua-grama.jpg}
        {grayscale_rua-grama.jpg}
        {segmented_rua-grama.jpg}
        {Rua e grama}
    \caption{Exemplos negativos e suas versões em grayscale e segmentadas.}
    \label{fig:exemplos-negativos}
\end{figure}
